\chapter{Reaction Theory II: Resonances and the \texorpdfstring{$R$}{R} Matrix}
\marginnoteinline{Source date: 01/08/2010.  Source pages 1 and 2 are duplicate scans; source page 11 is blank.}

\sourcepage{1}
\section{Single-level Breit--Wigner form}
The goal is to continue Reaction Theory I, derive the Breit--Wigner form, and begin the foundations for an elementary treatment of $R$-matrix theory.

From Reaction Theory I,
\begin{align*}
  f_\ell(k)
  &=\frac{1}{k}\ee^{\ii\delta_\ell}\sin\delta_\ell\\
  &=\frac{1}{k}\frac{\sin\delta_\ell}
  {\cos\delta_\ell-\ii\sin\delta_\ell}\\
  &=\frac{1}{k}\frac{1}{\cot\delta_\ell-\ii}.
  \tag{1}
\end{align*}
The partial-wave cross section is
\begin{equation*}
  \sigma_\ell=\frac{4\pi}{k^2}(2\ell+1)\sin^2\delta_\ell.
  \tag{2}
\end{equation*}
It is maximal when
\[
  \delta_\ell=\frac{(2n+1)\pi}{2},
\]
which corresponds to a resonance of the effective potential for the $\ell$th partial wave.

Let $E_0$ be the energy at which the resonance occurs.  Near $E_0$, $\cot\delta_\ell(E_0)=0$, and a Taylor expansion gives
\[
  \cot\delta_\ell(E)
  \simeq \cot\delta_\ell(E_0)
  +(E-E_0)\left[\frac{\dd}{\dd E}\cot\delta_\ell(E)\right]_{E=E_0}
  +\text{higher-order terms},
\]
or
\begin{equation*}
  \cot\delta_\ell(E)
  \simeq(E-E_0)\left[\frac{\dd}{\dd E}\cot\delta_\ell(E)\right]_{E=E_0}.
  \tag{3}
\end{equation*}

\sourcepage{2}
\textit{This source page duplicates source page 1.}

\sourcepage{3}
Define the resonance width
\begin{equation*}
  \Gamma_\ell
  \equiv-2\left[\frac{\dd}{\dd E}\cot\delta_\ell(E)\right]^{-1}_{E=E_0}.
  \tag{4}
\end{equation*}
Then
\[
  \cot\delta_\ell(E)\simeq-\frac{2}{\Gamma_\ell}(E-E_0).
\]
Substitution into the partial-wave amplitude gives
\begin{align*}
  f_\ell
  &=\frac{1}{k}\frac{1}{-2(E-E_0)/\Gamma_\ell-\ii}\\
  &=-\frac{1}{k}
  \frac{\Gamma_\ell/2}{E-E_0+\ii\Gamma_\ell/2}.
  \tag{5}
\end{align*}
Therefore
\begin{equation*}
  \sigma_\ell
  =4\pi(2\ell+1)|f_\ell|^2
  =\frac{4\pi}{k^2}(2\ell+1)
  \frac{\Gamma_\ell^2/4}
  {(E-E_0)^2+\Gamma_\ell^2/4}.
  \tag{6}
\end{equation*}
This is the Breit--Wigner formula for a resonant cross section.

\marginnoteinline{Question recorded in the notes: how do the cross sections factor into quantities that are actually obtained from, and used in, ENDF?  Additional self-notes point back to earlier sections and toward a more robust SLBW treatment, spinless one-channel reductions, and eventually two-channel/two-level examples.}

\sourcepage{4}
\section{Elastic and reaction cross sections in terms of the reflection factor}
Reaction Theory I gave
\begin{align*}
  \sigma_s
  &=\frac{4\pi}{k^2}\sum_{\ell=0}^{\infty}(2\ell+1)\sin^2\delta_\ell\\
  &=\pi\bar\lambda^2\sum_{\ell=0}^{\infty}(2\ell+1)|1-U_\ell|^2,
  \qquad \bar\lambda=\frac{1}{k},
  \tag{7}
\end{align*}
where the $\ell$-wave reflection factor is
\begin{equation*}
  U_\ell=\ee^{2\ii\delta_\ell}.
  \tag{8}
\end{equation*}
The partial elastic cross section is therefore
\begin{equation*}
  \sigma_{s,\ell}=\pi\bar\lambda^2(2\ell+1)|1-U_\ell|^2.
  \tag{9}
\end{equation*}

For reaction processes, let $N_{i,\ell}$ be the number of $\ell$-wave reactions per second per target nucleus:
\[
  N_{i,\ell}=\sigma_{r,\ell}\phi_{\mathrm{inc}},
\]
with the incident flux normalized to $v=\hbar k/m$.  Equivalently, $N_{i,\ell}$ is the net current through a large sphere about the target.  Inserting the appropriate asymptotic wave function into the current integral gives
\begin{equation*}
  \sigma_{r,\ell}=\pi\bar\lambda^2(2\ell+1)(1-|U_\ell|^2).
  \tag{11}
\end{equation*}
No reaction can occur without elastic scattering, as indicated by this decomposition.

\sourcepage{5}
Hence
\begin{align*}
  \sigma_{\mathrm{total},\ell}
  &=\sigma_{s,\ell}+\sigma_{r,\ell}\\
  &=\pi\bar\lambda^2(2\ell+1)
  \left[|1-U_\ell|^2+1-|U_\ell|^2\right]\\
  &=2\pi\bar\lambda^2(2\ell+1)
  \left[1-\operatorname{Re}(U_\ell)\right].
  \tag{12}
\end{align*}

\subsection*{Calculation of the reflection factor}
For $r\ge b$, write the reduced radial wave function as
\begin{equation*}
  w_\ell(r)=C_\ell\left[w_\ell^-(r)-U_\ell w_\ell^+(r)\right].
  \tag{13}\label{eq:rt2-external-wave}
\end{equation*}
The external functions are
\begin{equation*}
  w_\ell^\pm=kr\left[n_\ell(kr)\pm\ii j_\ell(kr)\right]
  =A_\ell(r)\pm\ii B_\ell(r),
  \tag{13b}
\end{equation*}
which are spherical Hankel functions in the notation of the notes.

Define the logarithmic derivative
\begin{equation*}
  D_\ell\equiv b\left[\frac{\dd w_\ell/\dd r}{w_\ell}\right]_{r=b}.
  \tag{14}\label{eq:rt2-log-derivative}
\end{equation*}
Substitution of \eqref{eq:rt2-external-wave} into \eqref{eq:rt2-log-derivative} and solution for $U_\ell$ gives
\begin{equation*}
  U_\ell
  =\left[
  \frac{D_\ell-b(w_\ell^{-\prime}/w_\ell^-)}
  {D_\ell-b(w_\ell^{+\prime}/w_\ell^+)}
  \frac{w_\ell^-}{w_\ell^+}
  \right]_{r=b}.
  \tag{15}
\end{equation*}

\sourcepage{6}
Further define
\begin{equation*}
  b\left[\frac{\dd w_\ell^\pm/\dd r}{w_\ell^\pm}\right]_{r=b}
  \equiv a_\ell\pm\ii b_\ell.
  \tag{16}
\end{equation*}
Using $w_\ell^\pm=A_\ell\pm\ii B_\ell$ gives the shift factor
\begin{equation*}
  a_\ell
  =b\left[
  \frac{A_\ell'A_\ell+B_\ell'B_\ell}
  {A_\ell^2+B_\ell^2}
  \right]_{r=b},
  \tag{17}
\end{equation*}
and the penetration factor
\begin{equation*}
  b_\ell
  =b\left[
  \frac{A_\ell'B_\ell-B_\ell'A_\ell}
  {A_\ell^2+B_\ell^2}
  \right]_{r=b}.
  \tag{18}
\end{equation*}
Both are tabulated functions.  Define the hard-sphere phase $\Delta_\ell$ by
\begin{equation*}
  \left[\frac{w_\ell^-}{w_\ell^+}\right]_{r=b}
  \equiv\ee^{2\ii\Delta_\ell}.
  \tag{19}
\end{equation*}
For an impenetrable sphere, the phase $\Delta_\ell$ is the full phase shift, and the reflection factor reduces to the ratio $w_\ell^-/w_\ell^+$, as expected.

\sourcepage{7}
\section{Elementary \texorpdfstring{$R$}{R}-matrix construction}
$R$-matrix theory can be viewed as an accurate and relatively convenient functional form to which experimental data in the resolved-resonance region can be fitted.  The fitted parameters $E_\lambda$ and $\gamma_{\lambda c}$ describe, in a suitable sense, resonance locations and widths.

The notes first consider simple spinless particles and a scalar partial wave, postponing the full matrix representation.  Expand the internal wave function as
\begin{equation*}
  \phi_\ell(E,r)=\sum_\lambda A_{\lambda\ell}\phi_\ell(E_\lambda,r).
  \tag{20}\label{eq:rt2-internal-expansion}
\end{equation*}
The functions obey
\begin{equation*}
  \left\{\frac{\dd^2}{\dd r^2}
  +\frac{2m}{\hbar^2}\left[E-V(r)
  -\frac{\ell(\ell+1)\hbar^2}{2mr^2}\right]\right\}
  \phi_\ell(E,r)=0,
  \tag{21}\label{eq:rt2-physical-equation}
\end{equation*}
\begin{equation*}
  \left\{\frac{\dd^2}{\dd r^2}
  +\frac{2m}{\hbar^2}\left[E_\lambda-V(r)
  -\frac{\ell(\ell+1)\hbar^2}{2mr^2}\right]\right\}
  \phi_\ell(E_\lambda,r)=0.
  \tag{22}\label{eq:rt2-basis-equation}
\end{equation*}
A common boundary condition is imposed at $r=b$:
\begin{equation*}
  \left[\frac{\dd\phi_\ell(E_\lambda,r)}{\dd r}\right]_{r=b}
  =b^{-1}B_\ell\phi_\ell(E_\lambda,b),
  \tag{23}\label{eq:rt2-basis-boundary}
\end{equation*}
with the same constant $B_\ell$ for the complete basis.

\sourcepage{8}
By orthogonality,
\begin{equation*}
  A_{\lambda\ell}
  =\int_0^b\phi_\ell(E_\lambda,r)\phi_\ell(E,r)\dd r.
  \tag{24}
\end{equation*}
Multiply \eqref{eq:rt2-physical-equation} and \eqref{eq:rt2-basis-equation} by the opposite wave functions and subtract.  The result can be arranged as
\begin{equation*}
  0=\frac{2m}{\hbar^2}(E-E_\lambda)
  \phi_\ell(E,r)\phi_\ell(E_\lambda,r)
  -\phi_\ell''(E,r)\phi_\ell(E_\lambda,r)
  +\phi_\ell''(E_\lambda,r)\phi_\ell(E,r).
  \tag{25}
\end{equation*}
Integrating from $0$ to $b$ gives
\begin{equation*}
  A_{\lambda\ell}
  =\frac{\hbar^2}{2m(E-E_\lambda)}
  \left[
  \phi_\ell'(E_\lambda,r)\phi_\ell(E,r)
  -\phi_\ell'(E,r)\phi_\ell(E_\lambda,r)
  \right]_{r=b}.
  \tag{26}
\end{equation*}
Use the boundary condition \eqref{eq:rt2-basis-boundary} in the expansion \eqref{eq:rt2-internal-expansion}.  At the boundary,
\begin{equation*}
  \phi_\ell(E,b)
  =\frac{\hbar^2}{2mb}
  \sum_\lambda
  \frac{\phi_\ell^2(E_\lambda,b)}{E_\lambda-E}
  \left[b\phi_\ell'(E,b)-B_\ell\phi_\ell(E,b)\right].
  \tag{27}\label{eq:rt2-boundary-expansion}
\end{equation*}

\sourcepage{9}
The $R$ matrix---a scalar $R$ function here---is therefore
\begin{align*}
  R_\ell(E)
  &=\frac{\hbar^2}{2mb}
  \sum_\lambda
  \frac{\phi_\ell^2(E_\lambda,b)}{E_\lambda-E}
  \tag{28a}\\
  &=\sum_\lambda\frac{\gamma_{\lambda\ell}^2}{E_\lambda-E},
  \tag{28b}
\end{align*}
where the reduced-width amplitude is
\begin{equation*}
  \gamma_{\lambda\ell}
  \equiv\sqrt{\frac{\hbar^2}{2mb}}\,\phi_\ell(E_\lambda,b).
  \tag{29}
\end{equation*}
From \eqref{eq:rt2-boundary-expansion},
\begin{equation*}
  R_\ell(E)
  =\left[
  \frac{\phi_\ell(E,r)}
  {r\phi_\ell'(E,r)-B_\ell\phi_\ell(E,r)}
  \right]_{r=b}.
  \tag{30}\label{eq:rt2-R-logderivative}
\end{equation*}
For $B_\ell=0$, this is the previously used reciprocal logarithmic derivative.

Using the exterior form
\[
  \phi_\ell(E,r)=C_\ell\left[\phi_\ell^-(E,r)-U_\ell\phi_\ell^+(E,r)\right],
\]
substitution into \eqref{eq:rt2-R-logderivative} gives a direct relation between $R_\ell$ and $U_\ell$.

\sourcepage{10}
Define the external logarithmic derivatives
\[
  L_\ell^\pm
  \equiv\left[
  \frac{r}{\phi_\ell^\pm}
  \frac{\dd\phi_\ell^\pm}{\dd r}
  \right]_{r=b}.
\]
Then the notes record
\begin{equation*}
  U_\ell
  =\left.\frac{\phi_\ell^-}{\phi_\ell^+}\right|_{r=b}
  \frac{1-R_\ell(L_\ell^- -B_\ell)}
  {1-R_\ell(L_\ell^+ -B_\ell)}.
  \tag{33}
\end{equation*}
This relates the internal $R$ function to the observable reflection factor.

The multichannel matrix form is sketched schematically as
\[
  \bm U
  =\bm\rho^{1/2}\bm\phi_{\mathrm{out}}
  [\bm I-\bm R(\bm L-\bm B)]^{-1}
  [\bm I-\bm R(\bm L^*-\bm B)]
  \bm\phi_{\mathrm{inc}}\bm\rho^{-1/2},
\]
with notation left for later development.  The $R$ matrix has elements $R_{cc'}$, where $c$ and $c'$ label entrance and exit channels.  A channel label collects the particle type, spins, and other quantum numbers needed to specify a particular partial wave.  The source ends with: ``Next time: level matrix and simplifications.''

\sourcepage{11}
\textit{Blank source page.}
